\documentclass[11pt,a4paper]{article}

% XeLaTeX/LuaLaTeX : fontspec + polyglossia ; repli pdflatex avec les mêmes polices.
\usepackage{iftex}
\ifPDFTeX
  \usepackage[T1]{fontenc}
  \usepackage[utf8]{inputenc}
  \usepackage{lmodern}
  \usepackage[french]{babel}
  \frenchsetup{StandardLists=true, AutoSpacePunctuation=false, og=«, fg=»}
  \addto\captionsfrench{\renewcommand{\figurename}{Fig.}\renewcommand{\tablename}{Tab.}}
\else
  \usepackage{fontspec}
  \usepackage{polyglossia}
  \setmainlanguage{french}
\fi
\usepackage[a4paper,margin=2.5cm]{geometry}
\usepackage{microtype}
\usepackage{amssymb}
\usepackage{booktabs,tabularx,array}
\usepackage{enumitem}
\usepackage{titlesec}
\usepackage{fancyhdr}
\usepackage{xcolor}
\usepackage{caption}
\usepackage{float}
\usepackage{tikz}
\usetikzlibrary{positioning,arrows.meta,fit,backgrounds,shapes.geometric,calc}
\usepackage[hidelinks]{hyperref}

\setlength{\parskip}{0.5em}
\setlength{\parindent}{0pt}
\setlist{itemsep=0.15em,topsep=0.3em}
\captionsetup{font=small,labelfont=bf}
\renewcommand{\arraystretch}{1.25}
\newcolumntype{L}{>{\raggedright\arraybackslash}X}
\newcolumntype{P}[1]{>{\raggedright\arraybackslash}p{#1}}
\setlength{\emergencystretch}{3em}

\titleformat{\section}{\large\bfseries}{\thesection.}{0.6em}{}
\titleformat{\subsection}{\normalsize\bfseries}{\thesubsection}{0.6em}{}

\pagestyle{fancy}
\fancyhf{}
\lhead{\small INPT -- Smart ICT}
\rhead{\small IA agentique pour le RIC en O-RAN}
\cfoot{\small\thepage}
\renewcommand{\headrulewidth}{0.4pt}

% repère d'une hypothèse ajoutée (voir section « Hypothèses ajoutées »)
\newcommand{\hyp}[1]{\textsuperscript{\textbf{[H#1]}}}
\newcommand{\tool}[1]{\texttt{#1}}

\tikzset{
  box/.style={draw, rounded corners=2pt, align=left, inner sep=5pt, font=\small},
  cbox/.style={draw, rounded corners=2pt, align=center, inner sep=5pt, font=\small},
  arr/.style={-{Stealth[length=2.2mm]}, thick},
  lbl/.style={font=\footnotesize\itshape, inner sep=2pt},
  zone/.style={draw, dashed, rounded corners=3pt, fill=gray!6},
  zlbl/.style={font=\footnotesize\itshape, anchor=north west},
}

\begin{document}

%----------------------------------------------------------------- en-tête
\thispagestyle{plain}
\begin{center}
  {\small Institut National des Postes et Télécommunications (INPT), Rabat\\
  Filière Smart ICT -- 3\textsuperscript{e} année}\\[1.2em]
  {\LARGE\bfseries IA agentique pour le RIC en O-RAN}\\[0.4em]
  {\large Proposition d'implémentation -- version révisée}\\[1.2em]
  \begin{tabular}{@{}l@{\hspace{3em}}l@{}}
    \textit{Réalisé par :} & \textit{Encadré par :}\\
    Souhail Bourhim & Pr. Mustapha Benjillali\\
    Ghita Boulmani & \\
  \end{tabular}\\[1em]
  {\small Octobre 2026}
\end{center}
\vspace{0.5em}
\hrule
\vspace{0.8em}

\textbf{Résumé.} Ce document présente ce que nous proposons de réaliser concrètement
dans le cadre du projet, au-delà de l'étude bibliographique. L'idée est de construire,
sur une plateforme 5G open source, un contrôleur dans lequel un agent basé sur un LLM
traduit l'objectif de l'opérateur en consignes pour un contrôleur par apprentissage par
renforcement (RL), avec une couche de vérification entre les deux. Nous comparerons
ensuite cette approche à des solutions plus simples sur les mêmes scénarios.

Le projet s'appuie sur les trois articles de référence : l'architecture hiérarchique
« le LLM guide, le RL exécute » de Bao et al. (LLM-hRIC), le problème d'allocation de
ressources entre slices étudié par Wu et al., et l'approche multi-agents DRL de
Rezazadeh et al., qui servira de référence côté RL. L'ensemble tourne sur une seule
machine Linux, sans matériel radio.

Cette version ajoute Hermes Agent, un agent open source existant, comme simple
\emph{interface opérateur}. Hermes transmet les intentions à une passerelle qui n'expose
que trois outils (\tool{submit\_intent}, \tool{get\_kpis}, \tool{get\_status}). Il ne remplace
ni LangGraph, ni le vérificateur déterministe, ni la xApp RL ; il n'a aucun accès aux
commandes E2 ni au partage des PRB, et il reste hors de la boucle temps réel. Ce n'est
pas une intégration O-RAN native. Les hypothèses ajoutées sont repérées dans le texte par
\hyp{1} à \hyp{8} et listées en section~\ref{sec:hyp}.

%----------------------------------------------------------------- 1
\section{Contexte et problématique}

Une cellule 5G dispose d'une capacité radio limitée, mesurée en blocs de ressources
physiques (PRB). Cette capacité est partagée entre plusieurs \emph{slices}, chacun ayant
des besoins différents (figure~\ref{fig:slices}). La demande varie au cours de la
journée, donc la répartition des PRB entre slices doit être ajustée en continu.

\begin{figure}[H]
\centering
\begin{tikzpicture}[node distance=4mm]
  \node[cbox, fill=gray!15, minimum width=5.2cm, minimum height=7mm] (e) {\textbf{Slice eMBB}};
  \node[cbox, fill=gray!15, minimum width=4cm, minimum height=7mm, right=2mm of e] (u) {\textbf{Slice URLLC}};
  \node[cbox, fill=gray!15, minimum width=3.6cm, minimum height=7mm, right=2mm of u] (b) {\textbf{Best effort}};
  \node[font=\small, above=1mm of u] {Capacité radio d'une cellule (PRB)};
  \node[box, below=5mm of e, text width=4.8cm] (e2) {Vidéo, téléchargements\\\footnotesize Débit élevé, grande part des PRB};
  \node[box, below=5mm of u, text width=3.6cm] (u2) {Robots, contrôle à distance\\\footnotesize Latence très faible, fiabilité};
  \node[box, below=5mm of b, text width=3.2cm] (b2) {Smartphones classiques\\\footnotesize Utilise le reste};
  \draw (e) -- (e2); \draw (u) -- (u2); \draw (b) -- (b2);
\end{tikzpicture}
\caption{Une cellule, trois slices aux besoins différents, une répartition à ajuster en permanence.}
\label{fig:slices}
\end{figure}

Dans l'architecture O-RAN, cette décision est prise dans le RAN Intelligent Controller
(RIC). Le RIC héberge des applications de contrôle : les xApps dans le Near-RT RIC
(boucles de 10~ms à 1~s) et les rApps dans le Non-RT RIC (au-delà d'une seconde).

Côté opérateur, l'objectif est de pouvoir exprimer une intention en langage naturel,
par exemple « garantir 10~Mb/s au slice URLLC entre 18h et 19h », et de laisser le
réseau s'en charger. La question à laquelle nous voulons répondre est donc la suivante :

\begin{quote}
\itshape Comment permettre à un opérateur de piloter le RAN en langage naturel, tout en
gardant des décisions rapides et en empêchant l'IA de prendre des actions dangereuses ?
\end{quote}

%----------------------------------------------------------------- 2
\section{Limites des approches existantes}

\begin{table}[H]
\centering\small
\begin{tabularx}{\textwidth}{@{}P{2.8cm}LLL@{}}
\toprule
Approche & Principe & Avantage & Limite\\
\midrule
Règles fixes & Répartition statique, modifiée manuellement & Simple, prévisible & Trop lente et rigide face à une demande variable\\
RL seul & Une politique entraînée associe les KPI à une répartition des PRB & Rapide (ms), efficace en optimisation & Ne connaît que l'objectif appris ; un nouvel objectif impose un nouvel entraînement\\
LLM seul & Le LLM lit l'objectif et fixe lui-même la répartition & Comprend le langage naturel, explique ses choix & Trop lent pour les boucles rapides ($\sim$1~s ou plus) ; peut proposer des actions erronées\\
\bottomrule
\end{tabularx}
\caption{Comparaison des approches existantes.}
\end{table}

Le terme « agent » a deux sens dans les articles étudiés. Un agent RL est une politique
entraînée qui associe un état à une action, sans raisonnement. Un système d'IA agentique
est un LLM qui planifie, appelle des outils, observe le résultat et se corrige. Notre
proposition combine les deux.

%----------------------------------------------------------------- 3
\section{Solution proposée}

Nous séparons les rôles selon l'échelle de temps et selon le niveau de confiance.
L'opérateur s'exprime dans Hermes Agent, qui sert d'interface (section~\ref{sec:hermes}).
Le service LangGraph, placé dans le Non-RT RIC, interprète l'intention et propose un
objectif. Un vérificateur déterministe contrôle les limites et la sécurité de cet objectif
avant toute modification de l'objectif du RL. La xApp RL, dans le Near-RT RIC, reste seule
responsable des décisions rapides d'allocation des PRB. Hermes et les LLM restent hors de
la boucle temps réel (figure~\ref{fig:archi}).

\begin{figure}[H]
\centering
\begin{tikzpicture}[x=1cm,y=1cm]
  % zones
  \draw[zone] (0,1.35) rectangle (15.45,-0.95);
  \node[zlbl] at (0.1,1.3) {Interface opérateur (hors RIC, hors boucle temps réel)};
  \draw[zone] (5.75,-1.95) rectangle (15.4,-4.4);
  \node[zlbl] at (5.85,-2.0) {Non-RT RIC / SMO};
  \draw[zone] (5.75,-5.15) rectangle (15.4,-7.45);
  \node[zlbl] at (8.1,-5.2) {Near-RT RIC (FlexRIC)};

  % ligne 0 : interface
  \node[box, text width=2.0cm] (op) at (1.35,0) {\textbf{Opérateur}\\\footnotesize langage naturel};
  \node[box, text width=3.0cm] (he) at (5.9,0) {\textbf{Hermes Agent}\\\footnotesize interface de chat,\\\footnotesize client MCP};
  \node[box, text width=4.6cm] (gw) at (12.7,0) {\textbf{Passerelle MCP / plugin}\\\footnotesize\tool{submit\_intent}, \tool{get\_kpis},\\\footnotesize\tool{get\_status}};

  % ligne 1 : Non-RT RIC
  \node[box, text width=4.0cm, fill=gray!12] (lg) at (12.9,-3.3) {\textbf{Service LangGraph}\\\footnotesize interprète l'intention,\\\footnotesize propose un objectif RL};
  \node[box, text width=3.2cm] (vf) at (7.7,-3.3) {\textbf{Vérificateur}\\\footnotesize déterministe : limites,\\\footnotesize sécurité, journal};
  \node[box, dashed, text width=2.9cm] (cf) at (2.6,-3.3) {\textbf{Confirmation opérateur}\\\footnotesize démo uniquement\\\footnotesize console séparée};

  % ligne 2 : Near-RT RIC
  \node[box, text width=3.4cm] (rl) at (7.7,-6.5) {\textbf{xApp RL (DDPG)}\\\footnotesize décide seule de la\\\footnotesize répartition des PRB};
  \node[box, text width=4.0cm] (kp) at (12.9,-6.5) {\textbf{xApp KPM}\\\footnotesize supervision : débit, PRB,\\\footnotesize latence par slice};
  \node[box, dashed, text width=4.3cm] (nb) at (2.65,-6.5) {\textbf{Aucun accès pour Hermes} aux commandes E2, au contrôle radio ni à la modification brute du partage des PRB};

  % RAN
  \node[box, text width=9.0cm] (ran) at (10.55,-9.3) {\textbf{RAN et c\oe ur 5G open source}\\\footnotesize gNB OAI + UE simulés (rfsim) + c\oe ur OAI CN5G ; slices eMBB, URLLC, best effort};

  % flèches
  \draw[arr] (op) -- node[lbl,above]{intention} (he);
  \draw[arr] (he) -- node[lbl,above]{MCP / plugin} (gw);
  \draw[arr] (11.5,0 |- gw.south) -- node[lbl,left]{\tool{submit\_intent}} (11.5,0 |- lg.north);
  \draw[arr, dashed] (14.0,0 |- lg.north) -- node[lbl,right,align=left]{\tool{get\_kpis,}\\\tool{get\_status}} (14.0,0 |- gw.south);
  \draw[arr] (lg.west) -- (vf.east);
  \draw[arr, <->] (vf.west) -- node[lbl,above]{démo} (cf.east);
  \draw[arr] (vf.south) -- node[lbl,right]{politique A1} (rl.north);
  \draw[arr] (kp.north) -- node[lbl,right]{KPI} (lg.south);
  \draw[arr] (kp.west) -- (rl.east);
  \draw[arr] (rl.south) -- node[lbl,right]{contrôle E2} (rl.south |- ran.north);
  \draw[arr] (kp.south |- ran.north) -- node[lbl,right]{rapports E2} (kp.south);
\end{tikzpicture}
\caption{Architecture révisée. Hermes Agent est une couche d'interface sans accès au RIC ;
LangGraph, le vérificateur déterministe et la xApp RL gardent leur rôle. Traits pleins :
demandes et actions ; trait discontinu : lecture seule.}
\label{fig:archi}
\end{figure}

Notre contribution porte sur cinq points :
\begin{enumerate}
  \item un contrôle hiérarchique LLM + RL, inspiré de Bao et al., appliqué au problème
        d'allocation entre slices de Wu et al. ; le LLM n'intervient jamais dans la boucle rapide ;
  \item une couche de sécurité déterministe qui vérifie chaque objectif proposé par le LLM
        et compte les actions bloquées ;
  \item une mise en \oe uvre sur une vraie pile open source (OpenAirInterface + FlexRIC),
        alors que la plupart des travaux du domaine restent en simulation ;
  \item une interface opérateur volontairement étroite (Hermes Agent derrière une passerelle
        à trois outils), évaluée séparément du contrôle ;
  \item une évaluation comparative de quatre contrôleurs sur les mêmes scénarios (section~\ref{sec:eval}).
\end{enumerate}

%----------------------------------------------------------------- 4
\section{Fonctionnement du service LangGraph}

Nous utilisons un seul agent, développé avec LangGraph et exposé comme un petit service
HTTP. Il suit une boucle planification -- action -- observation (approche ReAct), avec une
vérification systématique avant chaque action (figure~\ref{fig:boucle}).

\begin{figure}[H]
\centering
\begin{tikzpicture}[node distance=7mm and 7mm, every node/.style={font=\small}]
  \node[cbox, minimum width=2.4cm] (p) {Analyse\\\footnotesize intention $\to$ JSON};
  \node[cbox, minimum width=2.4cm, right=of p] (pl) {Planification\\\footnotesize choix de l'action};
  \node[cbox, minimum width=2.4cm, right=of pl] (a) {Action\\\footnotesize appel d'un outil};
  \node[cbox, minimum width=2.4cm, right=of a, fill=gray!15] (v) {Vérification\\\footnotesize règles fixes};
  \node[cbox, minimum width=2.4cm, below=12mm of v] (ap) {Application\\\footnotesize envoi au RIC};
  \node[cbox, minimum width=2.4cm, below=12mm of a] (o) {Observation\\\footnotesize lecture des KPI};
  \node[cbox, minimum width=2.4cm, below=12mm of pl] (r) {Évaluation\\\footnotesize objectif atteint ?};
  \node[cbox, minimum width=2.4cm, below=12mm of p, fill=gray!15] (f) {Fin\\\footnotesize mémorisation};
  \node[lbl, above=5mm of p, align=center] (in) {intention reçue\\de la passerelle};
  \draw[arr] (in) -- (p);
  \draw[arr] (p) -- (pl); \draw[arr] (pl) -- (a); \draw[arr] (a) -- (v);
  \draw[arr] (v) -- node[lbl,right,align=left]{valide\\(+ confirmation\\en démo)} (ap);
  \draw[arr] (ap) -- (o); \draw[arr] (o) -- (r);
  \draw[arr] (r) -- node[lbl,above]{oui} (f);
  \draw[arr] (r) -- node[lbl,right]{non} (pl);
  \draw[arr] (v.north) -- ++(0,5mm) -| node[lbl,above,pos=0.25]{action bloquée : replanifier} ([xshift=3mm]pl.north);
\end{tikzpicture}
\caption{Boucle de décision du service LangGraph.}
\label{fig:boucle}
\end{figure}

Le service n'agit jamais directement sur le réseau : il ne peut qu'appeler les outils
internes ci-dessous, et tout appel qui modifie la configuration passe par le vérificateur.
Ces outils internes ne sont pas ceux que voit Hermes (section~\ref{sec:hermes}).

\begin{table}[H]
\centering\small
\begin{tabularx}{\textwidth}{@{}P{5.6cm}L@{}}
\toprule
Outil interne & Rôle\\
\midrule
\tool{get\_kpis(slice, fenetre)} & Lit le débit, l'usage des PRB et la latence récents\\
\tool{set\_rl\_objective(poids, contraintes)} & Modifie l'objectif optimisé par la xApp RL (levier principal)\\
\tool{set\_slice\_share(parts)} & Fixe directement la répartition des PRB. Désactivé dans le contrôleur proposé ; utilisé seulement pour le contrôleur « LLM seul » de l'évaluation, toujours derrière le vérificateur\hyp{6}\\
\tool{wait\_and\_observe(secondes)} & Laisse le réseau se stabiliser avant la vérification suivante\\
\bottomrule
\end{tabularx}
\caption{Outils internes du service LangGraph.}
\end{table}

Pour le modèle, nous prévoyons un LLM local via Ollama (un modèle ouvert d'environ
8 milliards de paramètres avec appel d'outils, par exemple Qwen). Il est gratuit,
fonctionne hors ligne et reste réaliste pour un déploiement en périphérie du réseau.
Nous testerons aussi un modèle plus grand accessible par API (Claude ou GPT) sur les
mêmes scénarios, afin de mesurer le gain en précision, en latence et en coût.

Certains travaux récents (MX-AI, AgentRAN) utilisent plusieurs agents. Pour un projet
d'un semestre à deux, nous préférons un seul agent, plus simple à mettre au point. Si le
temps le permet, nous ajouterons un agent « analyste » en lecture seule, chargé
d'expliquer les anomalies observées dans les KPI.

%----------------------------------------------------------------- 5
\section{Interface opérateur : Hermes Agent}
\label{sec:hermes}

Pour que l'opérateur formule ses intentions dans une interface de conversation existante,
nous plaçons Hermes Agent en amont du service LangGraph. Hermes Agent est un agent open
source (licence MIT) développé par Nous Research, qui se connecte à des outils externes
par le protocole MCP ou par des plugins\hyp{1}. Ce n'est pas un composant O-RAN et il n'a
aucun support natif du RIC : du point de vue du réseau, c'est un client MCP ordinaire. Il
ne remplace ni LangGraph, ni le vérificateur déterministe, ni la xApp RL, et le rôle du
contrôleur RL ne change pas.

\subsection{Chaîne de traitement}

\begin{enumerate}
  \item L'opérateur formule une intention dans Hermes.
  \item Hermes la transmet à la passerelle (serveur MCP ou plugin Hermes personnalisé)
        avec l'outil \tool{submit\_intent}.
  \item La passerelle transmet l'intention au service LangGraph, qui l'interprète et propose un objectif\hyp{5}.
  \item Le vérificateur déterministe contrôle les limites et la sécurité de l'objectif
        proposé. Un objectif rejeté n'atteint jamais la xApp RL.
  \item Dans le scénario de démonstration, l'opérateur confirme dans une console séparée
        avant l'application\hyp{4}.
  \item L'objectif validé est transmis à la xApp RL par une politique A1. La xApp RL reste
        seule responsable des décisions rapides d'allocation des PRB dans le Near-RT RIC.
  \item L'opérateur suit l'avancement dans Hermes avec \tool{get\_status} et \tool{get\_kpis}.
\end{enumerate}

Hermes et les LLM ne font donc jamais partie de la boucle temps réel.

\subsection{Outils exposés par la passerelle}

La passerelle expose exactement trois outils, et aucun autre.

\begin{table}[H]
\centering\small
\begin{tabularx}{\textwidth}{@{}P{3.8cm}LL@{}}
\toprule
Outil & Rôle & Effet sur le réseau\\
\midrule
\tool{submit\_intent(texte)} & Transmet au service LangGraph le texte de l'intention, tel qu'écrit par l'opérateur\hyp{2} ; renvoie un identifiant et l'état initial & Crée une demande. Rien n'est appliqué sans validation du vérificateur (et confirmation en démo)\\
\tool{get\_kpis(slice, fenetre)} & Renvoie les KPI récents (débit, PRB, latence) & Lecture seule\\
\tool{get\_status(id)} & Renvoie l'état de l'intention (reçue, interprétée, rejetée avec motif, en attente de confirmation, appliquée, expirée) et l'objectif RL en vigueur & Lecture seule\\
\bottomrule
\end{tabularx}
\caption{Outils exposés à Hermes par la passerelle.}
\end{table}

\subsection{Ce que Hermes ne peut pas faire}

\begin{itemize}
  \item Hermes ne reçoit \textbf{aucun accès direct} aux commandes E2, au contrôle radio ni à
        la modification brute du partage des PRB.
  \item Aucun outil exposé ne modifie la configuration : \tool{submit\_intent} ne fait que
        soumettre une demande.
  \item Toute action qui modifie la configuration passe par le service de vérification.
  \item La confirmation opérateur ne passe pas par Hermes : sinon le LLM pourrait se
        l'accorder lui-même\hyp{4}.
\end{itemize}

La configuration MCP de Hermes permet de n'enregistrer, pour chaque serveur, qu'une liste
blanche d'outils (clé \tool{include}) ; nous y déclarons les trois outils de la passerelle.
Ces garanties ne reposent pas seulement sur la liste d'outils. D'après sa documentation,
Hermes est un agent généraliste qui peut aussi disposer d'outils comme un terminal, des
fichiers ou le web\hyp{3}. Pour ce projet, nous désactivons ces outils et nous exécutons
Hermes dans un conteneur dédié dont le seul accès réseau est le port de la passerelle. La
passerelle n'a elle-même qu'un interlocuteur, le service LangGraph, sans route vers FlexRIC
ni vers OAI\hyp{5}.

\subsection{Prototype minimal}

Le prototype montre le circuit complet avec un périmètre réduit. Il est d'abord réalisé
contre le simulateur (le service LangGraph s'appuie alors sur l'environnement gym), puis
contre la plateforme OAI + FlexRIC si elle est prête\hyp{8}.

\begin{enumerate}
  \item Un serveur MCP en Python implémente les trois outils et appelle l'API HTTP du service LangGraph.
  \item Hermes est configuré avec ce serveur comme seule source d'outils.
  \item Scénario de démonstration : l'opérateur écrit « garantir 10~Mb/s au slice URLLC
        pendant 10 minutes ». Hermes appelle \tool{submit\_intent} ; la passerelle renvoie
        l'état « en attente de confirmation » ; l'opérateur confirme dans la console
        séparée ; l'objectif est appliqué ; l'opérateur demande où en est le réseau et Hermes
        affiche l'état et les KPI par slice avec \tool{get\_status} et \tool{get\_kpis}.
  \item Cas de rejet : une intention hors limites (par exemple un débit irréaliste) est
        refusée par le vérificateur, et le motif apparaît dans \tool{get\_status}.
  \item Test négatif : une demande hors périmètre (« fixe le partage des PRB à 90/5/5 »)
        ne doit produire aucun changement d'état, car aucun outil ne le permet.
\end{enumerate}

Le prototype est réussi si ces trois cas se déroulent comme prévu et si le journal de la
passerelle ne contient que les trois outils exposés.

%----------------------------------------------------------------- 6
\section{Plateforme expérimentale}

La plateforme tourne sur un ordinateur portable (Ryzen~9, 16 c\oe urs, 32~Go de RAM,
GPU RTX~5070 de 8~Go) sous Ubuntu dans WSL2, avec Docker. Le lien radio est simulé,
aucun matériel SDR n'est nécessaire. Le module SCTP, nécessaire aux interfaces du c\oe ur
OAI, se charge correctement dans WSL2.

\begin{table}[H]
\centering\small
\begin{tabularx}{\textwidth}{@{}P{3.6cm}LP{2.6cm}@{}}
\toprule
Composant & Outil & Réutilisé / développé\\
\midrule
Station de base et UE & gNB et nrUE OpenAirInterface, mode rfsim, 2 à 3 UE & Réutilisé\\
C\oe ur 5G & OAI CN5G avec 2 à 3 slices & Réutilisé\\
Near-RT RIC & FlexRIC & Réutilisé\\
Trafic & iperf3 sur chaque UE, charge variable & Scénarios à écrire\\
xApp de supervision & Modèle de service KPM de FlexRIC + stockage des mesures & Développé\\
xApp de contrôle & Modèle de service de contrôle des slices de FlexRIC & Développé\\
Politique RL & DDPG, entraîné dans un simulateur Python (gym) puis déployé & Développé\\
Service LangGraph et vérificateur & LangGraph, Ollama, API HTTP & Développé\\
Interface opérateur & Hermes Agent, dans un conteneur isolé, outils généralistes désactivés\hyp{1}\hyp{3} & Réutilisé, configuré\\
Passerelle & Serveur MCP en Python (ou plugin Hermes), trois outils & Développé\\
Console de confirmation & Ligne de commande ou petite page web du service de vérification\hyp{4} & Développé (minimal)\\
\bottomrule
\end{tabularx}
\caption{Composants de la plateforme.}
\end{table}

Nous associons FlexRIC à OAI plutôt qu'à srsRAN, car le contrôle des slices dans
FlexRIC est mieux supporté avec OAI ; c'est aussi la combinaison utilisée dans MX-AI.
L'entraînement du RL se fait d'abord en simulation, car il serait trop lent sur la
plateforme émulée.

%----------------------------------------------------------------- 7
\section{Méthode d'évaluation}
\label{sec:eval}

Nous menons deux évaluations séparées. La première mesure le contrôle du réseau ; la
seconde mesure l'interface. Ainsi, la performance de l'interface n'est pas confondue avec
celle du contrôle RL, ni l'inverse.

\subsection{Évaluation du contrôle : quatre contrôleurs}

Les mêmes scénarios de trafic sont exécutés avec quatre contrôleurs. Hermes n'intervient
dans aucun d'eux : les intentions sont injectées directement dans le service LangGraph, et
seules les intentions correctement extraites (vérifiées au préalable) sont retenues pour
cette comparaison\hyp{7}.

\begin{table}[H]
\centering\small
\begin{tabularx}{\textwidth}{@{}P{5cm}L@{}}
\toprule
Contrôleur & Ce qu'il représente\\
\midrule
Répartition statique égale & Référence sans intelligence\\
RL seul & Approche de type Rezazadeh et al. / Wu et al.\\
LLM seul, action directe & Approche agentique naïve\\
\textbf{LLM (objectifs) + RL (exécution)} & \textbf{Notre proposition}, d'après Bao et al.\\
\bottomrule
\end{tabularx}
\caption{Contrôleurs comparés.}
\end{table}

\textbf{Scénarios}, tous générés avec iperf3 pour que chaque contrôleur voie le même trafic :
\begin{enumerate}
  \item charge stable, puis pic soudain sur un slice ;
  \item nouvelle intention en cours d'exécution (« protéger URLLC pendant 10 minutes ») ;
  \item deux intentions contradictoires.
\end{enumerate}

\textbf{Métriques :}
\begin{itemize}
  \item proportion du temps où l'intention de chaque slice est respectée ;
  \item débit et latence par slice ;
  \item temps de réaction après un pic de charge ou une nouvelle intention ;
  \item nombre d'actions dangereuses proposées puis bloquées par le vérificateur ;
  \item latence et coût du LLM du service LangGraph (modèle local vs modèle distant).
\end{itemize}

\subsection{Évaluation de l'interface : Hermes et passerelle}

Cette évaluation s'arrête à la décision du vérificateur : elle n'utilise ni le débit ni la
latence radio, et ses résultats ne servent pas à classer les contrôleurs. Elle reprend le
scénario des intentions ambiguës ou mal formulées, auparavant dans la comparaison des
contrôleurs\hyp{7}.

\begin{table}[H]
\centering\small
\begin{tabularx}{\textwidth}{@{}P{2.9cm}LL@{}}
\toprule
Mesure & Définition & Source des données\\
\midrule
Latence de l'interface & Temps entre l'envoi du message de l'opérateur et l'arrivée de \tool{submit\_intent} à la passerelle, puis entre la réponse de la passerelle et son affichage ; médiane et 95\textsuperscript{e} centile sur des répétitions. Le temps de décision du service LangGraph est exclu & Horodatages de la passerelle et journal de Hermes\\
Fidélité de transmission & Part des intentions arrivées à \tool{submit\_intent} sans reformulation par le LLM de Hermes\hyp{2} & Texte saisi comparé au texte reçu par la passerelle\\
Qualité d'extraction de l'intention & Sur un jeu d'intentions annotées à la main\hyp{8} : exactitude par champ (slice, métrique, valeur cible, fenêtre de temps), part des intentions entièrement correctes, détection des cas ambigus & Sortie structurée du service LangGraph comparée aux annotations\\
Appels d'outils & Par outil : appels réussis, rejetés par la passerelle (arguments invalides), rejetés par le vérificateur (hors limites), en erreur & Journaux de la passerelle et du vérificateur\\
Tentatives hors périmètre & Demandes d'une action non exposée (par exemple fixer directement le partage des PRB) : aucun changement d'état attendu & Journal et état du RIC avant et après\\
\bottomrule
\end{tabularx}
\caption{Mesures de l'interface, indépendantes du contrôle RL.}
\end{table}

%----------------------------------------------------------------- 8
\section{Plan de travail}

Le travail est organisé en cinq phases et deux axes menés en parallèle, avec un point de
validation à la fin de chaque phase (figure~\ref{fig:plan}). L'axe B n'attend pas la
plateforme : le simulateur permet de commencer le RL et l'agent dès la première phase. La
phase 4 (interface) ne démarre que lorsque le service LangGraph et son vérificateur
fonctionnent ; elle peut se réduire à un client en ligne de commande si le temps manque
(section~\ref{sec:risques}). La répartition des axes entre nous deux reste à fixer.

\begin{figure}[H]
\centering
\begin{tikzpicture}[x=1cm,y=1cm]
  \node[font=\small\bfseries] at (6.0,0.5) {Axe A : plateforme, xApps};
  \node[font=\small\bfseries] at (12.2,0.5) {Axe B : RL, agent, interface};

  \node[align=left, anchor=west, font=\small] at (0,-0.7) {\textbf{Phase 1}\\\footnotesize Fondations};
  \node[cbox, text width=5.5cm, minimum height=1.1cm] at (6.0,-0.7) {\footnotesize Installer OAI, c\oe ur, FlexRIC ;\\\footnotesize xApp de supervision (KPM)};
  \node[cbox, text width=5.5cm, minimum height=1.1cm] at (12.2,-0.7) {\footnotesize Simulateur gym des 3 slices};
  \node[font=\footnotesize\itshape, anchor=west] at (3.1,-1.575) {$\blacklozenge$ KPI en direct};

  \node[align=left, anchor=west, font=\small] at (0,-2.45) {\textbf{Phase 2}\\\footnotesize Contrôle};
  \node[cbox, text width=5.5cm, minimum height=1.1cm] at (6.0,-2.45) {\footnotesize xApp de contrôle des slices};
  \node[cbox, text width=5.5cm, minimum height=1.1cm] at (12.2,-2.45) {\footnotesize Entraînement DDPG\\\footnotesize en simulation};
  \node[font=\footnotesize\itshape, anchor=west] at (3.1,-3.325) {$\blacklozenge$ Le RL modifie la répartition des PRB};

  \node[align=left, anchor=west, font=\small] at (0,-4.2) {\textbf{Phase 3}\\\footnotesize Agent};
  \node[cbox, text width=5.5cm, minimum height=1.1cm] at (6.0,-4.2) {\footnotesize Relier l'agent au RIC\\\footnotesize (politique A1, données KPM)};
  \node[cbox, text width=5.5cm, minimum height=1.1cm] at (12.2,-4.2) {\footnotesize Service LangGraph, vérificateur,\\\footnotesize console de confirmation};
  \node[font=\footnotesize\itshape, anchor=west] at (3.1,-5.075) {$\blacklozenge$ Intention satisfaite de bout en bout (sans Hermes)};

  \node[align=left, anchor=west, font=\small] at (0,-5.95) {\textbf{Phase 4}\\\footnotesize Interface};
  \node[cbox, text width=5.5cm, minimum height=1.1cm] at (6.0,-5.95) {\footnotesize Passerelle MCP à 3 outils ;\\\footnotesize tests de refus (hors périmètre)};
  \node[cbox, text width=5.5cm, minimum height=1.1cm] at (12.2,-5.95) {\footnotesize Configurer Hermes ; intentions\\\footnotesize annotées ; mesures d'interface};
  \node[font=\footnotesize\itshape, anchor=west] at (3.1,-6.825) {$\blacklozenge$ Hermes soumet une intention et affiche l'état et les KPI};

  \node[align=left, anchor=west, font=\small] at (0,-7.7) {\textbf{Phase 5}\\\footnotesize Résultats};
  \node[cbox, fill=gray!12, text width=11.6cm, minimum height=1.1cm] at (9.1,-7.7) {\footnotesize Évaluation du contrôle (4 contrôleurs) et de l'interface ;\\\footnotesize analyse, rapport, préparation de la soutenance (en commun)};
\end{tikzpicture}
\caption{Phases du projet. Les losanges marquent les points de validation.}
\label{fig:plan}
\end{figure}

%----------------------------------------------------------------- 9
\section{Risques et solutions de repli}
\label{sec:risques}

\begin{table}[H]
\centering\small
\begin{tabularx}{\textwidth}{@{}LL@{}}
\toprule
Risque & Solution de repli\\
\midrule
Le contrôle des slices OAI + FlexRIC ne fonctionne pas de façon fiable & Garder la supervision KPM sur la plateforme réelle et faire tourner la boucle de contrôle en simulation (ns-O-RAN ou notre simulateur)\\
Machine trop limitée pour la plateforme et un LLM local & Utiliser le LLM via API ou sur une seconde machine\\
La politique RL entraînée en simulation se comporte mal sur la plateforme & Court ré-entraînement sur la plateforme, ou exécutant à règles simples sous le LLM\\
LLM trop lent ou peu fiable dans l'appel d'outils & Prompts plus courts, schéma JSON imposé, ou modèle plus grand\\
Hermes difficile à installer ou instable (version, environnement) & La même passerelle est pilotée par un petit client en ligne de commande ; le contrôle et son évaluation n'en dépendent pas\\
Le LLM de Hermes appelle mal les outils ou reformule l'intention & Texte transmis tel quel, validation de schéma à la passerelle, modèle plus grand ; mesuré dans l'évaluation de l'interface\\
Hermes contourne la passerelle avec ses outils généralistes & Outils désactivés, conteneur isolé sans route vers FlexRIC ni OAI, tests négatifs du prototype\\
Dérive du périmètre (mémoire, compétences, messageries de Hermes) & N'utiliser que le chat et MCP ; limiter la phase 4 dans le temps ; abandon possible au profit du client en ligne de commande\\
Deux LLM en chaîne ajoutent de la latence et des erreurs & Interface hors boucle temps réel ; latence mesurée à part ; bascule vers le client direct si nécessaire\\
\bottomrule
\end{tabularx}
\caption{Principaux risques identifiés.}
\end{table}

%----------------------------------------------------------------- 10
\section{Hypothèses ajoutées dans cette version}
\label{sec:hyp}

Les points suivants ont été ajoutés ou choisis pour intégrer Hermes Agent. Ils ne viennent
pas de la version précédente et restent à valider avec l'encadrant ou par des essais.

\begin{description}[leftmargin=2.6em,labelwidth=2.2em,labelsep=0.4em,font=\bfseries]
  \item[H1] Hermes Agent (Nous Research, open source, licence MIT) sert uniquement
        d'interface de conversation et de client MCP. D'après sa documentation, il prend en
        charge les serveurs MCP et les plugins ; nous n'avons pas encore testé son
        installation sur notre machine (WSL2) et la version à utiliser reste à figer.
  \item[H2] Le modèle utilisé par Hermes est configurable ; nous supposons qu'il peut s'agir
        du modèle local (Ollama) ou d'un modèle distant, ce qui reste à vérifier. En local, Hermes
        et le service LangGraph partageraient le même modèle servi par Ollama, car le GPU
        (8~Go, dont une partie utilisée par Windows) ne peut pas en charger deux. Hermes
        transmet le texte de l'opérateur tel quel à \tool{submit\_intent} ; l'extraction
        structurée est faite par LangGraph, ce qui permet de la mesurer à part.
  \item[H3] Hermes est isolé : conteneur dédié, outils généralistes désactivés, accès réseau
        limité à la passerelle. La possibilité de désactiver ces outils par configuration
        est à confirmer dans la version utilisée.
  \item[H4] La confirmation opérateur se fait dans une console séparée du service de
        vérification (ligne de commande ou petite page web), pas par un outil Hermes. Une
        demande non confirmée expire après un délai configurable (valeur à fixer). Les
        exécutions d'évaluation automatisées n'utilisent pas la confirmation ; seule la
        démonstration l'exige.
  \item[H5] La passerelle n'a qu'un interlocuteur, l'API HTTP du service LangGraph ;
        \tool{get\_kpis} et \tool{get\_status} y sont servis en lecture seule (les KPI viennent
        du stockage de la xApp de supervision).
  \item[H6] L'outil interne \tool{set\_slice\_share} n'est plus un moyen de secours du
        contrôleur proposé : il est désactivé pour celui-ci et sert seulement au contrôleur
        « LLM seul », derrière le vérificateur. Il n'est jamais exposé à Hermes.
  \item[H7] Pour comparer les quatre contrôleurs, les intentions sont injectées sans Hermes et
        seules les intentions correctement extraites sont retenues. Le scénario des intentions
        ambiguës passe dans l'évaluation de l'interface.
  \item[H8] Le jeu d'évaluation de l'extraction compte de l'ordre de 40 intentions annotées à
        la main (français et anglais, dont des cas ambigus) ; ce dimensionnement est une
        estimation à ajuster. Les mesures d'interface se font d'abord sur le simulateur, puis
        sur la plateforme si elle est prête.
\end{description}

%----------------------------------------------------------------- refs
\section*{Références}
\small
\begin{enumerate}[label={[\arabic*]}]
  \item Rezazadeh et al., apprentissage par renforcement profond multi-agents pour le slicing du RAN.
  \item Wu et al., IA agentique pilotée par LLM pour la résilience en O-RAN.
  \item Bao et al., LLM-hRIC : contrôle hiérarchique LLM + RL dans le RIC.
  \item MX-AI, agents LLM sur OAI + FlexRIC, arXiv:2508.09197.
  \item AgentRAN, hiérarchie d'agents auto-organisée pour l'Open RAN, arXiv:2508.17778.
  \item xTRUCE, arbitrage sûr entre xApps pilotées par LLM, arXiv:2608.28532.
  \item OpenAirInterface (\url{https://openairinterface.org}), FlexRIC (\url{https://gitlab.eurecom.fr/mosaic5g/flexric}), LangGraph, Ollama.
  \item Hermes Agent, Nous Research, open source, licence MIT (\url{https://hermes-agent.nousresearch.com} ; dépôt \url{https://github.com/NousResearch/hermes-agent} ;
        configuration MCP : \url{https://hermes-agent.nousresearch.com/docs/user-guide/features/mcp}).
  \item Model Context Protocol (MCP) (\url{https://modelcontextprotocol.io}).
\end{enumerate}

\end{document}
